% !TEX program = lualatex
% Transparencias de Fundamentos de las Matemáticas II
% Clase 02: Números enteros II

\documentclass[aspectratio=169,11pt]{beamer}

\usepackage{fontspec}
\usepackage[spanish,es-nodecimaldot]{babel}
\usepackage{amsmath,amssymb,mathtools}
\usepackage{booktabs}
\usepackage{csquotes}
\usepackage{xcolor}

\usetheme{Madrid}
\usecolortheme{default}
\setbeamertemplate{navigation symbols}{}
\setbeamertemplate{footline}[frame number]

\definecolor{FdMBlue}{RGB}{20,55,100}
\definecolor{FdMRed}{RGB}{130,30,30}
\definecolor{FdMGreen}{RGB}{25,100,60}

\setbeamercolor{structure}{fg=FdMBlue}
\setbeamercolor{title}{fg=white,bg=FdMBlue}
\setbeamercolor{frametitle}{fg=white,bg=FdMBlue}
\setbeamercolor{block title}{fg=white,bg=FdMBlue}
\setbeamercolor{block body}{bg=blue!3}
\setbeamercolor{alerted text}{fg=FdMRed}

\setmainfont{Latin Modern Roman}
\setsansfont{Latin Modern Sans}
\setmonofont{Latin Modern Mono}

\newcommand{\N}{\mathbb{N}}
\newcommand{\Z}{\mathbb{Z}}
\newcommand{\Q}{\mathbb{Q}}
\newcommand{\R}{\mathbb{R}}
\newcommand{\C}{\mathbb{C}}
\newcommand{\Pow}{\mathcal{P}}
\newcommand{\id}{\operatorname{id}}
\newcommand{\mcd}{\operatorname{mcd}}
\newcommand{\mcm}{\operatorname{mcm}}
\newcommand{\im}{\operatorname{Im}}

\title[FdM II -- Clase 02]{Números enteros II}
\subtitle{Valor absoluto y división euclídea}
\author{Antonio Falcó Montesinos}
\institute{Universidad CEU Cardenal Herrera}
\date{Fundamentos de las Matemáticas II}

\begin{document}

\begin{frame}
  \titlepage
\end{frame}

\begin{frame}{Objetivos de la clase}
\begin{itemize}
  \item Definir valor absoluto en \(\mathbb{Z}\).
  \item Interpretar el valor absoluto como distancia al cero.
  \item Introducir la división euclídea.
\end{itemize}
\end{frame}

\begin{frame}{Mapa de la clase}
\tableofcontents
\end{frame}

\section{Valor absoluto}

\begin{frame}{Valor absoluto}
\begin{block}{Definición}
\begin{itemize}
  \item Si \(a\geq 0\), entonces \(|a|=a\).
  \item Si \(a<0\), entonces \(|a|=-a\).
  \item El valor absoluto siempre es un entero no negativo.
\end{itemize}
\end{block}
\end{frame}

\begin{frame}{Valor absoluto}
\begin{block}{Interpretación}
\begin{itemize}
  \item \(|a|\) mide la distancia de \(a\) al cero.
  \item Por eso \(|a|=|-a|\).
  \item También permite expresar distancia entre enteros como \(|a-b|\).
\end{itemize}
\end{block}
\end{frame}

\begin{frame}{Valor absoluto}
\begin{block}{Propiedades}
\begin{itemize}
  \item \(|a|=0\) si y solo si \(a=0\).
  \item \(|ab|=|a||b|\).
  \item \(|a+b|\leq |a|+|b|\).
\end{itemize}
\end{block}
\end{frame}

\section{División euclídea}

\begin{frame}{División euclídea}
\begin{block}{Enunciado}
\begin{itemize}
  \item Si \(a,b\in\mathbb{Z}\) y \(b>0\), existen enteros únicos \(q,r\) tales que \(a=bq+r\).
  \item El resto satisface \(0\leq r<b\).
  \item \(q\) es el cociente y \(r\) el resto.
\end{itemize}
\end{block}
\end{frame}

\begin{frame}{División euclídea}
\begin{block}{Ejemplo}
\begin{itemize}
  \item Para \(a=37\) y \(b=5\): \(37=5\cdot 7+2\).
  \item El cociente es \(7\).
  \item El resto es \(2\).
\end{itemize}
\end{block}
\end{frame}

\begin{frame}{División euclídea}
\begin{block}{Cuidado}
\begin{itemize}
  \item El resto no puede ser negativo.
  \item El resto debe ser menor que el divisor.
  \item La unicidad depende de imponer \(0\leq r<b\).
\end{itemize}
\end{block}
\end{frame}

\section{Profundización}

\begin{frame}{Profundización}
\begin{block}{Dividendo negativo}
\begin{itemize}
  \item Para \(-17\) entre \(5\), buscamos \(-17=5q+r\) con \(0\leq r<5\).
  \item Tomando \(q=-4\), queda \(-17=5(-4)+3\).
  \item El resto es \(3\), no \(-2\).
\end{itemize}
\end{block}
\end{frame}

\begin{frame}{Profundización}
\begin{block}{Uso posterior}
\begin{itemize}
  \item La división euclídea fundamenta el algoritmo de Euclides.
  \item Permite definir congruencias mediante restos.
  \item Es central en divisibilidad y aritmética modular.
\end{itemize}
\end{block}
\end{frame}

\begin{frame}{Profundización}
\begin{block}{Comprobación rápida}
\begin{itemize}
  \item ¿Cumple el resto \(0\leq r<b\)?
  \item ¿Se recupera el dividendo al calcular \(bq+r\)?
  \item ¿El divisor se ha tomado positivo?
\end{itemize}
\end{block}
\end{frame}


\section{Detalle y práctica}

\begin{frame}{Detalle y práctica}
\begin{block}{Valor absoluto y distancia}
\begin{itemize}
  \item El valor absoluto mide distancia al cero: \(|a|\geq 0\).
  \item La distancia entre dos enteros se escribe \(|a-b|\).
  \item La desigualdad triangular resume una propiedad geométrica: \(|a+b|\leq |a|+|b|\).
\end{itemize}
\end{block}
\end{frame}

\begin{frame}{Detalle y práctica}
\begin{block}{Actividad breve}
\begin{itemize}
  \item Calcular \(|-8|\), \(|3-11|\) y \(|-4+9|\).
  \item Encontrar todos los \(x\in\mathbb{Z}\) tales que \(|x|\leq 3\).
  \item Explicar por qué \(|a|=0\) implica \(a=0\).
\end{itemize}
\end{block}
\end{frame}

\begin{frame}{Cierre}
\begin{itemize}
  \item El valor absoluto expresa distancia y tamaño sin signo.
  \item La división euclídea produce cociente y resto únicos.
  \item El resto siempre se elige entre \(0\) y \(b-1\).
\end{itemize}
\end{frame}

\begin{frame}{Trabajo recomendado}
\begin{itemize}
  \item Releer las secciones correspondientes del manual.
  \item Identificar definiciones, símbolos nuevos y enunciados principales.
  \item Preparar dudas y ejemplos para el seminario asociado.
\end{itemize}
\end{frame}

\end{document}
